\begin{statementbox}

\textbf{\textcolor{CStatement}{它要解决什么问题？}}

有些分式在某个位置附近，分子越来越接近 $0$，分母也越来越接近 $0$。
此时不能把 $\frac00$ 当作答案：两者接近 $0$ 的快慢不同，最终比值也可能不同。
洛必达法则提供一种计算办法：\textbf{在条件满足时，把分子、分母分别求导，再比较求导后的结果。}

\medskip
\begin{center}\begin{tikzpicture}[font=\small,>=stealth,line join=round]
\draw[gray!45,rounded corners=3pt] (0,0) rectangle (14.2,9.4);
\node at (7.1,8.8) {同样趋于零，比值仍可能是不同的数};
\begin{scope}[shift={(3.8,4.7)},scale=1.3]
\draw[->,gray] (-1.3,0)--(1.45,0) node[above] {$t=x-1$};
\draw[->,gray] (0,-2.7)--(0,2.9) node[left] {$y$};
\draw[blue!75!black,thick,samples=70,domain=-.8:.8] plot (\x,{2*\x+\x*\x});
\draw[orange!85!black,thick,dashed] (-.8,-.8)--(.8,.8);
\draw[violet,thick,samples=70,domain=-.8:.8] plot (\x,{\x*\x});
\draw[green!45!black,thick,dashdotted] (-.8,-2.4)--(.8,2.4);
\fill (0,0) circle(2pt);
\end{scope}
\node[anchor=west,align=left] at (8.4,5.2)
{$\textcolor{orange!85!black}{\text{分母 }g=t}$\\[15pt]
$\textcolor{blue!75!black}{(2t+t^2)/t\to2}$\\[15pt]
$\textcolor{violet}{t^2/t\to0}$\\[15pt]
$\textcolor{green!45!black}{3t/t\to3}$};
\node[text width=13.2cm,align=center] at (7.1,.65) {四条曲线都经过原点；“都趋于零”没有决定它们高度之比。};
\end{tikzpicture}
\end{center}
\end{statementbox}
\clearpage
\begin{statementbox}
\textbf{\textcolor{CStatement}{主要结论：分子、分母都趋近于 $0$}}


设 $x$ 越来越接近 $a$（但 $x\ne a$），且
\[
\lim_{x\to a}f(x)=0,\qquad \lim_{x\to a}g(x)=0.
\]
还需要满足：
\begin{enumerate}[leftmargin=*, itemsep=0.25em, topsep=0.25em]
    \item 在 $a$ 附近、但不取 $a$ 的地方，$f(x)$ 和 $g(x)$ 都能求导，且 $g'(x)\ne0$；
    \item $\displaystyle\lim_{x\to a}\frac{f'(x)}{g'(x)}$ 能确定为一个数，或确定地趋向 $+\infty$、$-\infty$。
\end{enumerate}
那么
\[
\boxed{\displaystyle
\lim_{x\to a}\frac{f(x)}{g(x)}
=\lim_{x\to a}\frac{f'(x)}{g'(x)}.}
\]
这里的 $f'(x)$、$g'(x)$ 表示\textbf{分别对分子、分母求导}，不是对整个分式求导。

\medskip
\textbf{\textcolor{CStatement}{另一个可用情况}}

如果分子、分母的大小都越来越大、没有上界，满足相应的求导条件，并且求导后的比值有确定的极限，也可用同样的办法。
例如 $x\to+\infty$ 时，$\frac{2x^2+x}{x^2+1}$ 的分子、分母都越来越大。

\medskip
\textbf{\textcolor{CStatement}{高中学习提醒}}

这个名称所指的一般结论，完整证明超出了通常的高中课堂推导范围。
下面只用高中范围内的导数定义、因式分解、三角关系，\textbf{严格解释几类常见情况}；
不把这些例子的证明说成整个法则的完整证明。遇到学校要求的常规解法时，应优先使用因式分解、有理化等方法。

\end{statementbox}
